\documentclass[11pt]{article}
\usepackage[margin=1in]{geometry}
\usepackage[hidelinks]{hyperref}
\usepackage{booktabs}
\usepackage{amsmath}
\usepackage{amssymb}
\usepackage{enumitem}
\usepackage{xcolor}

\title{Differentiable RDycore, Step One:\\
Adjoint-Based Manning Roughness Calibration\\
\large A proposed first paper and development plan}
\author{Mark Adams (LBNL) \\ \small prepared for Gautam Bisht and the RDycore team}
\date{August 18, 2026}

\begin{document}
\maketitle

\section{Summary}

I propose we build a \emph{discrete adjoint} capability into RDycore using
PETSc's \texttt{TSAdjoint}~\cite{zhang2022tsadjoint} and drive it to one
concrete, publishable endpoint: an end-to-end system that reads surface
observations (USGS stream gauges first), runs RDycore forward and adjoint
solves inside a gradient-based optimization loop (PETSc/TAO), and outputs a
\textbf{calibrated Manning roughness map} on the Houston 1\,km Hurricane
Harvey mesh. The adjoint returns the gradient of a gauge-misfit functional
with respect to \emph{every} cell's roughness at the cost of roughly one
extra forward solve, independent of the number of parameters --- this is
what makes distributed calibration tractable where trial-and-error and
ensemble screening are not.

This is deliberately the smallest slice of the larger differentiable-RDycore
plan (adjoint $\rightarrow$ calibration $\rightarrow$ PyTorch coupling
$\rightarrow$ learned physics) that yields a standalone paper and a strong
exhibit for next year's program recompete.

\section*{Status addendum (August 18, evening): implemented and verified}
\addcontentsline{toc}{section}{Status addendum}

\fbox{\parbox{0.97\linewidth}{\textbf{The plan below has been executed.}
On branch \texttt{adams/manning-draft}, increments 0, 1, 1b, 2a--2d, 3,
4, 5, and 6 (twin) are implemented, finite-difference-verified, and
committed, with 12 new CI tests (all green; full log:
\texttt{plans/RESULTS-manning-draft.md}; paper draft:
\texttt{papers/manning-calibration/}). Later sections are left in plan
tense for reference.}}

\medskip
\begin{center}\small
\begin{tabular}{@{}lll@{}}
\toprule
Verified result & Value & Gate \\
\midrule
Exact Roe flux Jacobian vs FD (wet jump / transcritical) & $7.6\times10^{-9}$ / $2.4\times10^{-10}$ & $10^{-6}$/$10^{-5}$ \\
Assembled matrix vs FD coloring, full matrix, 3 BC combos & $1.5$--$1.9\times10^{-8}$ & $10^{-6}$ \\
Adjoint $dJ/du_0$ and $dJ/dn$ vs FD (RK4, BEULER, ARKIMEX) & $\sim10^{-8}$ / $\sim3\times10^{-6}$ & $10^{-5}$ \\
Two-zone Manning recovery (TAO/BLMVM twin) & $4.1\times10^{-7}$ & $2\times10^{-2}$ \\
Per-cell recovery, 20 obs times + Tikhonov & $8.0\%$ & --- \\
Houston Harvey map (2{,}746 cells, implicit, 300-it prelim.) & $26.3\%$, map written & --- \\
\bottomrule
\end{tabular}
\end{center}

Highlights beyond the plan: the flux Jacobian is \emph{exact} (a
hand-written forward-mode differential of the flux code, entropy-fix
branches included); \textbf{implicit stepping now works} --- backward
Euler and ARK-IMEX (friction implicit, per-cell solves), both with
FD-gated adjoints --- and was \emph{required}: explicit friction NaNs at
Houston's $\Delta t = 30$\,s exactly as Section~\ref{sec:imex}
predicts. Discovered along the way: a new explicit source treatment was
prerequisite (the existing splittings embed $\Delta t$ in the RHS);
PETSc's TSEULER has no adjoint; the ARKIMEX adjoint needs a zero
explicit-part parameter Jacobian registered (else SEGV). Remaining:
converged Houston run (in progress), real-gauge calibration --- gauge
and datum data are ready (\texttt{data/harvey\_gauges/}), blocked only
on identifying the Houston mesh's projected CRS (team question).

\section{What exists and what is missing}

The RDycore SWE solver today is a forward model: explicit finite-volume
shallow water (Roe fluxes; HLLC planned but not yet implemented) with no
Jacobian, no adjoint, and no gradient path to its parameters. There is a worked in-tree precedent to
learn from --- the heat development branch (not a baseline for this
effort) carries an implicit time integrator, a Jacobian matrix, and
\texttt{TSSetIJacobian} with analytic per-cell derivatives, including a
libCEED Q-function Jacobian --- but those Jacobians are all
\emph{diagonal} (local source terms). What is missing is precisely this
plan's contribution: the \emph{coupled}, off-diagonal flux Jacobian on
the cell-adjacency graph, and all adjoint/sensitivity machinery.
Everything needed is either stock PETSc, demonstrated in-tree, or
mechanical to add:

\begin{itemize}[nosep]
  \item Manning $n$ is already a per-cell material property with per-region
        and per-domain public setters --- the parameter infrastructure exists.
  \item The heat branch provides worked reference examples for Jacobian
        matrix creation, implicit-TS and IJacobian registration, and a
        CEED Q-function Jacobian; the SWE drag-source Jacobian is the
        structural twin of its source Jacobian ($3\times3$ diagonal
        blocks instead of scalars).
  \item The SWE right-hand-side Jacobian is block-sparse ($3\times3$ blocks
        on the cell-adjacency graph); Roe flux Jacobians are closed-form.
        The friction-source Jacobian is diagonal, and
        $\partial S/\partial n$ has just two nonzeros per cell.
  \item \texttt{TSAdjoint} differentiates PETSc's own time integrators
        (including our explicit RK) with optimal trajectory checkpointing;
        we supply only the two Jacobian callbacks and a cost integrand.
  \item Gauge-site plumbing is native: the YAML
        \texttt{observations.sites} convention and its natural-ordering
        scatter (\texttt{src/time\_series.c}) already select cells to
        observe; the calibration work adds only observed-\emph{values}
        input and the misfit.
  \item An ARK-IMEX path is half-built: \texttt{TEMPORAL\_ARK\_IMEX}
        and a ``friction on the LHS'' source mode exist in the config
        schema and in the CEED source kernels --- only the integrator
        wiring and the source Jacobian are missing
        (Section~\ref{sec:imex}).
  \item The Houston 1\,km Harvey case is in the repository, and the
        USGS Harvey gauge data is downloaded (746 stage/discharge
        gauges with locations, HydroShare DOI
        10.4211/hs.c037167e497546a1bc1508dfb32a9cff, CC-BY; the
        remaining QC item is the stage-datum conversion). The LETKF
        twin-experiment driver provides the observation plumbing and
        the driver pattern to copy.
\end{itemize}

The work is CPU-backend only for the SWE operator --- the heat branch's
example shows CEED can express \emph{pointwise/diagonal} Jacobians as
Q-functions (a viable later route for a GPU drag-source Jacobian), but
there is no CEED path to the coupled off-diagonal flux-Jacobian assembly
the adjoint needs, so the CEED SWE kernels remain forward-only and
untouched. It lands behind a config flag with zero
impact on existing runs, and needs no new dependencies. Known technical risk: adjoints of SWE are
noisy at wetting/drying fronts; we validate on fully wet cases first and
treat wet/dry observations as future work.

\section{Related work: this is proven ground, and timely}

\paragraph{Adjoint Manning inversion is mature.}
Identification of distributed Manning coefficients by adjoint analysis goes
back two decades~\cite{ding2004manning}, with twin-experiment validation on
unstructured finite-volume SWE~\cite{yan2014adjoint} and GPU-accelerated
adjoint optimizers for 2D SWE~\cite{lacasta2016gpu}. The method itself is
not research risk.

\paragraph{Two production-grade analogs exist --- neither on our terrain.}
DassFlow (INSA Toulouse) is a 2D finite-volume shallow-water code built for
variational data assimilation, inferring Manning, bathymetry, and discharge,
including from SWOT satellite altimetry~\cite{brisset2018swot,dassflow}.
Thetis (Firedrake) calibrated a spatially varying Manning field for the
entire Baltic Sea with an automatically generated adjoint%
~\cite{karna2023baltic}, after storm-surge sensitivity studies that
carefully treated the overfitting problem for distributed
friction~\cite{warder2021storm} and an extension to coupled
morphodynamics~\cite{clare2022morpho}. Both validate our approach ---
and neither runs on DOE exascale machines or couples to an Earth system
model.

\paragraph{The differentiable-solver wave is cresting on exactly our use case.}
A 2026 paper, Inunda~\cite{li2026inunda}, rewrote a simplified
(local-inertial, raster) flood solver entirely in PyTorch and calibrated
roughness by gradient descent against gauges --- validated on a
\emph{Hurricane Harvey hindcast in Harris County} (0.67\,m MAE against
high-water marks). Hydrograd~\cite{liu2025uswe} (Julia) couples a 2D SWE
solver with neural flow-resistance terms; AegirJAX~\cite{cardoso2026aegir}
does differentiable coastal hydrodynamics in JAX; CaMa-Flood-GPU%
~\cite{kang2026cama} names gradient-based Manning calibration as its goal;
and differentiable rainfall--runoff models already infer Manning $n$ with
embedded neural networks~\cite{bindas2024routing}.

\paragraph{Our differentiator.}
Every code in that last group \emph{rewrote} a reduced solver inside an ML
framework to obtain differentiability; the FE-framework codes cannot reach
DOE leadership machines. RDycore + \texttt{TSAdjoint} is, to my knowledge,
the only path that makes an existing, validated (471M cells), GPU-capable,
E3SM-coupled DOE production code differentiable \emph{without rewriting it}.
Inunda's Harvey result is the perfect foil: our Phase-3 demo is the same
named U.S. flood event, in a code that also runs coupled continental-scale
simulations on Frontier.

\section{Incremental plan to the Manning map}
\label{sec:plan}

Each increment is small, independently testable (CTest against
finite-difference ground truth throughout), and lands behind a default-off
config flag. Detailed version in
\texttt{plans/manning-map-incremental-plan.md}.

\begin{center}
\begin{tabular}{@{}clp{5.6cm}p{4.6cm}@{}}
\toprule
\# & Increment & Deliverable & RDycore-team input \\
\midrule
0 & Config + FD baseline & \texttt{TSSetRHSJacobian} wired with PETSc
    FD-coloring Jacobian; flag default off & YAML schema approval \\
1 & Source Jacobians & Analytic drag/bed-slope diagonal blocks,
    preallocated bs=3 matrix & Physics review ($h_{\min}$ floor,
    semi-implicit friction) \\
2a--d & Roe flux Jacobian & Split by risk: edge-level derivative
    harness; physical-flux + frozen-dissipation Jacobian (validates
    assembly on smooth flows); exact dissipation derivatives (the one
    hard item, with the harness as oracle and the frozen Jacobian as
    schedule offramp); BC blocks & Entropy-fix review;
    \textbf{Roe vs HLLC decision} \\
3 & State adjoint & Trajectory + gauge-misfit cost; $dJ/du_0$ vs FD;
    new adjoint driver & Observation conventions, gauge file format \\
4 & Parameter gradient & $dJ/dn$ per region vs FD; first
    \emph{sensitivity map} & \textbf{Public API review} (E3SM-facing) \\
5 & TAO loop (per region) & Twin experiment recovers two-zone $n$ &
    Bounds, test-case choice \\
6 & Per-cell $n$ + regularization & \textbf{Calibrated Manning map} on
    Houston1km; optional real gauges & \textbf{Harvey gauge data};
    regularization prior; acceptance criteria \\
7 & Paper & Methods + twin + Harvey results & Co-authors, venue, scope \\
\bottomrule
\end{tabular}
\end{center}

Rough total: $\sim$1500 lines of new C (one library source file, one
driver, a small API surface, five CTests). Increments 4--6 each produce a
citable artifact (sensitivity map $\rightarrow$ recovered zones
$\rightarrow$ calibrated map).

\section{Strategic decisions --- application-science lead}
\label{sec:decisions}

These are the choices that shape the paper and the recompete story, and
they belong to the RDycore/application side, not the solver side:

\begin{enumerate}[nosep]
  \item \textbf{Observation strategy for paper 1.} Synthetic-only twin
        experiment vs.\ real USGS Harvey stage data (my recommendation:
        twin for validation, real gauges for the headline result if the
        data QC is manageable). The literature notes velocity observations
        constrain $n$ more strongly than stage but converge less
        readily~\cite{yan2014adjoint} --- do we have or want any velocity
        data? Satellite: SWOT water-surface elevation is proven usable in
        this exact inversion~\cite{brisset2018swot} and is a natural second
        paper; flood-extent imagery interacts badly with wet/dry
        nondifferentiability and should stay future work.
  \item \textbf{Roughness parameterization and prior.} Per-region
        (well-posed, few parameters) vs.\ per-cell with regularization
        (the interesting map, but overfitting is the known
        failure mode~\cite{warder2021storm}). Choice of Tikhonov vs.\
        total-variation, and whether the prior $n_0$ comes from land-cover
        data --- that prior is application knowledge.
  \item \textbf{Flux function on the adjoint path.} Roe is the
        committed path (the only Riemann solver implemented; its
        entropy-fix kinks are handled by wet-case validation and
        one-sided derivatives). HLLC is the \emph{backup}, triggered
        only if gradient validation or the Harvey calibration shows the
        kinks poisoning gradients --- and it means implementing the
        forward HLLC flux first (\texttt{RIEMANN\_HLLC} is enum-only
        today), then its Jacobian. Physics team owns the trigger call.
  \item \textbf{Target case and wet/dry exposure.} Houston1km Harvey is
        the compelling demo but exercises wetting/drying; validation
        cases stay fully wet. How wet is the calibration window we pick?
  \item \textbf{Public API shape.} What enters \texttt{rdycore.h}
        (\texttt{RDySetObservations}, \texttt{RDyAdjointSolve},
        \texttt{RDyGetSensitivities}, \ldots) and whether Fortran/E3SM
        bindings are needed now or later.
  \item \textbf{Venue and positioning.} WRR, JAMES, or GMD; and whether we
        engage Inunda~\cite{li2026inunda} head-on with the same Harvey
        event (direct comparison, higher impact, more work) or cite and
        differentiate.
  \item \textbf{Recompete framing.} How much of the PyTorch-bridge/AI-DA
        arc appears in paper 1's outlook section versus being held for the
        proposal narrative.
\end{enumerate}

\section{Future work (enabled, not included)}
\label{sec:future}

The calibrated-map system is the substrate for: zero-copy PyTorch coupling
via the existing PETSc$\leftrightarrow$PyTorch bridge (RDycore as a
\texttt{torch.autograd.Function}); neural source terms trained through the
solve (model-error correction, surface--groundwater exchange); hybrid
ensemble--variational DA with the LETKF line; SWOT and velocity
observations; second-order adjoints for uncertainty on the recovered map;
optimal gauge placement; implicit/IMEX stepping of the SWE system ---
IMEX with implicit friction is a near-term easy win quantified in
Section~\ref{sec:imex}, and fully implicit stepping on the new coupled
Jacobian follows; and sediment-parameter estimation. Each is a proposal work package with the
first paper as its demonstrated foundation.

\section{Appendix: equations, algorithms, and the PETSc API they map to}
\label{sec:math}

\subsection*{Forward model (exists today)}

RDycore solves the 2D shallow water equations in conservative variables
$u = (h, q_x, q_y)$ with $\mathbf{q} = (q_x,q_y) = h\mathbf{v}$:
\begin{align}
  \partial_t h + \nabla\!\cdot\!\mathbf{q} &= Q_r(x,t),
  \label{eq:mass}\\
  \partial_t \mathbf{q}
    + \nabla\!\cdot\!\Big(\tfrac{\mathbf{q}\otimes\mathbf{q}}{h}
    + \tfrac{g h^2}{2}\,\mathbb{I}\Big)
  &= -\,g h \nabla z_b \;-\; g\, n^2 h^{-7/3}\,\mathbf{q}\,\lVert\mathbf{q}\rVert,
  \label{eq:mom}
\end{align}
with rain/runoff forcing $Q_r$, bed elevation $z_b$, and Manning drag
written in momentum form (equivalently $-C_D\lVert\mathbf{v}\rVert\mathbf{v}$
with $C_D = g n^2 h^{-1/3}$), exactly as coded in
\texttt{src/swe/swe\_petsc.c}. Cell-centered finite volumes on an
unstructured DMPlex mesh give the semidiscrete system
\begin{equation}
  \frac{du_i}{dt}
  = -\frac{1}{|\Omega_i|}\sum_{e \in \partial\Omega_i}
      \mathcal{F}\big(u_i, u_j, \mathbf{n}_e\big)\,L_e
    \;+\; S\big(u_i;\, n_i, t\big)
  \;\equiv\; f_i(u, p),
  \label{eq:semidiscrete}
\end{equation}
where $\mathcal{F}$ is the Roe approximate Riemann flux (the only
Riemann solver implemented today; an upwinded variant exists, HLLC does
not), with well-balancing and an $h_{\min}$ wet/dry floor. Time
integration is explicit: forward Euler or explicit Runge--Kutta.

\emph{PETSc realization:} mesh and state via \texttt{DMPlex}; the RHS
$f(u,p)$ registered with
\texttt{TSSetRHSFunction(ts, rhs, OperatorRHSFunction, rdy)}
(\texttt{src/rdysetup.c}); integrator via \texttt{TSSetType}
(\texttt{TSEULER}/\texttt{TSRK}); windowed advance via
\texttt{TSSetMaxTime}/\texttt{TSSolve}.

\subsection*{State Jacobian (increments 0--2)}

Each interior edge $e=(i,j)$ contributes two $3\times3$ blocks per cell
row, so the RHS Jacobian is block-sparse on the cell-adjacency graph:
\begin{equation}
  J = \frac{\partial f}{\partial u}, \qquad
  J_{ii} \mathrel{+}= -\frac{L_e}{|\Omega_i|}\,
     \frac{\partial \mathcal{F}}{\partial u_L}
     + \frac{\partial S_i}{\partial u_i}, \qquad
  J_{ij} = -\frac{L_e}{|\Omega_i|}\,
     \frac{\partial \mathcal{F}}{\partial u_R},
  \label{eq:jac}
\end{equation}
with closed-form Roe flux Jacobians $\partial\mathcal{F}/\partial u_{L,R}$
and the local drag/bed-slope block
$\partial S/\partial u$ (from Eq.~\ref{eq:mom}; note
$\partial(h^{-7/3})/\partial h$ terms are regularized by the same
$h_{\min}$ floor as the forward model). Boundary conditions modify
diagonal blocks only.

\emph{PETSc realization:} matrix from \texttt{DMCreateMatrix} with
adjacency preallocation; blocked insertion via
\texttt{MatSetValuesBlocked} (bs$=3$, \texttt{MATBAIJ} on CPU,
\texttt{MATAIJ} for device compatibility); registered with
\texttt{TSSetRHSJacobian}; verified against finite differences with
\texttt{SNESComputeJacobianDefaultColor} (\texttt{MatFDColoring}).

\subsection*{Cost functional, trajectory, and discrete adjoint (increments 3--4)}

The gauge misfit over observation times $t_k$ with sparse observation
operator $H$ (one row per gauge, selecting $h$ at the gauge cell) is
\begin{equation}
  J(p) = \tfrac12 \sum_k
    \big(H u(t_k) - y_k\big)^{\!\top} R^{-1} \big(H u(t_k) - y_k\big),
  \label{eq:cost}
\end{equation}
where $y_k$ is the vector of observed stages at time $t_k$ and $R$ is
the observation-error covariance --- diagonal in practice,
$R = \mathrm{diag}(\sigma_g^2)$ with $\sigma_g$ the error standard
deviation of gauge $g$, so the misfit is the Gaussian negative
log-likelihood and the calibrated $n$ a maximum-likelihood estimate.
The cost is optionally written as an integral $\int_0^T r(u,t)\,dt$ for
the quadrature formulation.

In data-assimilation terms, Eqs.~\eqref{eq:cost} and \eqref{eq:opt}
together constitute \emph{strong-constraint 4D-Var} with the control
vector dominated by parameters rather than the initial state (the
regularization term of Eq.~\eqref{eq:opt} playing the role of the
background term): observations distributed over a window, the nonlinear
model as a hard constraint, and the gradient from one adjoint sweep ---
the same formulation DassFlow~\cite{dassflow} calls variational data
assimilation. This is also the precise sense in which the adjoint line
\emph{complements} the LETKF line: both minimize the same Gaussian
misfit, with ensembles as the natural minimizer for sequential state
correction and adjoints for whole-window parameter estimation, which is
what makes the hybrid EnVar step in Section~\ref{sec:future}
well-posed. \texttt{TSAdjoint} differentiates the
\emph{discrete} time stepper, so the backward sweep
\begin{equation}
  \dot\lambda = -\Big(\frac{\partial f}{\partial u}\Big)^{\!\top}\lambda
                - \Big(\frac{\partial r}{\partial u}\Big)^{\!\top},
  \qquad
  \dot\mu = -\Big(\frac{\partial f}{\partial p}\Big)^{\!\top}\lambda,
  \label{eq:adjoint}
\end{equation}
(run in its RK-discrete form) returns $dJ/du_0 = \lambda(0)$ and
$dJ/dp = \mu(0)$ exact to roundoff for the computed trajectory --- the
right property for FD validation and optimization. The parameter
Jacobian for Manning $n$ is two nonzeros per cell, from
Eq.~\eqref{eq:mom}:
\begin{equation}
  \frac{\partial S_{\mathbf{q},i}}{\partial n_i}
  = -\,2\, g\, n_i\, h_i^{-7/3}\, \mathbf{q}_i \lVert\mathbf{q}_i\rVert .
  \label{eq:dSdn}
\end{equation}

\emph{PETSc realization:} checkpointing via
\texttt{TSSetSaveTrajectory} (\texttt{TSTrajectory}; memory at test
scale, disk at Harvey scale); cost integrand via
\texttt{TSCreateQuadratureTS}; Eq.~\eqref{eq:dSdn} via
\texttt{TSSetRHSJacobianP}; sensitivities registered with
\texttt{TSSetCostGradients($\lambda$,$\mu$)} and computed by
\texttt{TSAdjointSolve}; $H$ as a sparse \texttt{MATAIJ}
(\texttt{MatCreateAIJ}), with gauge cells specified through the
existing \texttt{observations.sites} YAML convention
(\texttt{RDyObservationSites}) and the LETKF driver's observation
matrix as the pattern.

\subsection*{Bound-constrained calibration (increments 5--6)}

\begin{equation}
  \min_{\,n_{\mathrm{lo}} \le n \le n_{\mathrm{hi}}}\;
  J(n) \;+\; \tfrac{\beta}{2}\,\lVert L\,(n - n_0)\rVert^2,
  \label{eq:opt}
\end{equation}
per-region $n$ first (no regularization needed), then per-cell with
$L = \mathbb{I}$ (Tikhonov) or a discrete-gradient $L$
(smoothing/TV-like), prior $n_0$ from land cover.

\emph{PETSc realization:} \texttt{TaoCreate} with \texttt{TAOBLMVM} (or
\texttt{TAOBNCG}); bounds via \texttt{TaoSetVariableBounds}; the
objective/gradient callback runs one \texttt{TSSolve} +
\texttt{TSAdjointSolve} pair per iteration (gradient cost $\approx$ one
forward solve, independent of the number of parameters);
\texttt{TaoSolve} drives the loop.

\subsection*{Summary map}

\begin{center}
\begin{tabular}{@{}lllc@{}}
\toprule
Component & Equations & PETSc API & Increment \\
\midrule
Forward RHS $f(u,p)$ & \eqref{eq:mass}--\eqref{eq:semidiscrete} &
  \texttt{TSSetRHSFunction}, \texttt{DMPlex} & exists \\
Explicit stepping & --- & \texttt{TSSetType(TSEULER/TSRK)}, \texttt{TSSolve} & exists \\
State Jacobian $\partial f/\partial u$ & \eqref{eq:jac} &
  \texttt{TSSetRHSJacobian}, \texttt{MatSetValuesBlocked} & 0--2 \\
IMEX friction (Sec.~\ref{sec:imex}) & \eqref{eq:taufric} &
  \texttt{TSARKIMEX}, \texttt{TSSetIFunction}/\texttt{IJacobian} & 1b \\
FD verification & --- & \texttt{SNESComputeJacobianDefaultColor} & 1--2 \\
Trajectory checkpointing & --- & \texttt{TSSetSaveTrajectory} & 3 \\
Cost functional & \eqref{eq:cost} & \texttt{TSCreateQuadratureTS} & 3 \\
Observation operator $H$ & \eqref{eq:cost} & \texttt{MatCreateAIJ} & 3 \\
Adjoint sweep $\lambda,\mu$ & \eqref{eq:adjoint} &
  \texttt{TSSetCostGradients}, \texttt{TSAdjointSolve} & 3--4 \\
Parameter Jacobian $\partial f/\partial n$ & \eqref{eq:dSdn} &
  \texttt{TSSetRHSJacobianP} & 4 \\
Bounded optimization & \eqref{eq:opt} &
  \texttt{TAOBLMVM}, \texttt{TaoSetVariableBounds}, \texttt{TaoSolve} & 5--6 \\
\bottomrule
\end{tabular}
\end{center}

\section{Implicit time stepping: the IMEX-friction opportunity}
\label{sec:imex}

The Phase-1 source Jacobian enables a cheap, high-value integrator
upgrade that is worth calling out separately: \texttt{TSARKIMEX} with
the Manning drag term of Eq.~\eqref{eq:mom} treated implicitly and the
hyperbolic fluxes kept explicit. Because the drag is per-cell local, the
implicit stage matrix $I - a\,\Delta t\,\partial S/\partial u$ is
\emph{block-diagonal} with $3\times3$ blocks (effectively $2\times2$:
drag touches only the momentum rows and that cell's own $h$). The
``solve'' is an independent few-iteration Newton in every cell --- no
global system, no KSP, nothing to precondition, and trivially GPU
parallel. It requires only the increment-1 source Jacobian, none of the
flux blocks, and it replaces the current hand-rolled semi-implicit
friction splittings (\texttt{SEMI\_IMPLICIT} and the Xia--Liang-style
\texttt{IMPLICIT\_XQ2018}) with a principled, second-order, adaptive
(\texttt{TSAdapt}), and adjoint-compatible integrator.

Notably, this path is \emph{already half-built} on \texttt{main}:
\texttt{TEMPORAL\_ARK\_IMEX} and \texttt{SOURCE\_ARK\_IMEX} (``bed
friction terms moved to LHS'') exist in the config schema with
validation, and the CEED source kernels already branch on the IMEX
source mode --- but no \texttt{TSARKIMEX} is ever created and no
IFunction/IJacobian is registered. Completing it is exactly the
increment-1 source Jacobian plus integrator wiring. A related fact
strengthens the case: RDycore's existing adaptive stepping
(\texttt{target\_courant\_number} in the Harvey adaptive config)
targets the CFL limit \emph{only} --- in friction-bound thin-water
regimes that controller is flying blind, and implicit friction is what
makes aggressive CFL-targeted adaptivity safe.

\subsection*{How much larger a step? A stability estimate}

Explicit stepping faces two limits. The gravity-wave CFL condition,
which IMEX-friction retains,
\begin{equation}
  \Delta t_{\mathrm{CFL}} \;\approx\; \frac{C\,\Delta x}{\sqrt{gh} + |v|},
  \label{eq:cfl}
\end{equation}
and the friction-stiffness limit, which it removes: linearizing the
drag as $\dot{\mathbf q} = -\lambda\,\mathbf q$ with
$\lambda = g n^2 h^{-4/3}|v|$ gives explicit stability only for
\begin{equation}
  \Delta t \;\lesssim\; 2\tau_f, \qquad
  \tau_f = \frac{h^{4/3}}{g\,n^2\,|v|}.
  \label{eq:taufric}
\end{equation}
The IMEX step-size gain in a cell is the ratio
$\Delta t_{\mathrm{CFL}}/2\tau_f$; setting it to one gives a
\emph{binding-depth threshold}: friction is the limiter wherever
\begin{equation}
  h^{4/3} \;\lesssim\; \frac{\Delta x\, g\, n^2\, |v|}
                            {2C\,(\sqrt{gh}+|v|)} .
  \label{eq:threshold}
\end{equation}
For the Houston 1\,km mesh ($n=0.015$ in the current config,
$|v|\sim0.1$\,m/s) the threshold is $h \approx 5$\,cm --- and a
rainfall-driven Harvey domain spends much of the event with large areas
below it. Since the global explicit $\Delta t$ is the minimum over all
cells, a single thin wet cell drags the whole domain down to seconds;
with implicit friction the domain runs at the CFL of its
deepest/fastest cells.

\begin{center}
\begin{tabular}{@{}lllll@{}}
\toprule
Regime & $h$, $|v|$ & Friction limit $2\tau_f$ &
  CFL at $\Delta x{=}1$\,km & IMEX gain \\
\midrule
Dam break, post-break (ex2b) & 5\,m, 3\,m/s & $\sim$40\,min & CFL-bound
  & $1\times$ \\
Bayou in flood & 5\,m, 2\,m/s & $\sim$20\,min & $\sim$110\,s & $1\times$ \\
Sheet flow & 2\,cm, 0.1\,m/s & $\sim$35\,s & $\sim$2000\,s & $\sim50\times$ \\
Thin film / recession & 5\,mm, 0.1\,m/s & $\sim$8\,s & $\sim$4000\,s &
  $\sim500\times$ \\
Drying front & 1\,mm, 0.05\,m/s & $\sim$2\,s & $\sim$10$^4$\,s &
  $\sim10^3\times$ (capped by $h_{\min}$) \\
\bottomrule
\end{tabular}
\end{center}

Two cases anchor the estimate. In the \emph{deep} dam-break test (ex2b:
10\,m over 5\,m, $n=0.015$), $\tau_f \approx 20$\,minutes while the CFL
step is a fraction of a second --- friction is four orders of magnitude
from binding, so IMEX buys \emph{robustness, not speed}, in that phase.
On the Houston Harvey case the picture inverts after peak flood: the
long recession/drainage tail that dominates wall-clock is exactly where
the drag is stiff, and the achievable stable step improves by roughly
\textbf{one to two orders of magnitude} (approaching three at drying
fronts, where the $h_{\min}$ floor caps the estimate).

\subsection*{Why not just keep the semi-implicit splitting? The adjoint
argument}

For the \emph{forward} model, the existing splittings already bank most
of the step-size gain in the table, so the raw-speed argument for IMEX
is modest. The decisive argument is structural: the semi-implicit
treatments embed $\Delta t$ \emph{inside the RHS callback} (they
reconstruct the end-of-step momentum algebraically within $f$), so what
the integrator sees is not an ODE right-hand side --- it is consistent
only with forward Euler, and it admits no well-defined
$\partial f/\partial u$ to give to \texttt{TSSetRHSJacobian}.
\texttt{TSAdjoint}'s guarantee --- exact derivatives of the discrete
stepper --- assumes a genuine $f$. That leaves exactly two clean
options on the adjoint path: (a)~plain \emph{explicit} friction, which
is fine for the wet validation cases (increments 0--5) but collapses
$\Delta t$ to seconds in the thin-water phases of a Harvey window,
inflating adjoint cost and trajectory storage by the table's
$10$--$100\times$; or (b)~friction moved into \texttt{TSSetIFunction}
with its block-diagonal \texttt{TSSetIJacobian} under
\texttt{TSARKIMEX}, which \texttt{TSAdjoint} differentiates natively.
\textbf{We adopt (b)}: increment 1b completes the started
\texttt{SOURCE\_ARK\_IMEX} path and must land before the increment-6
Harvey calibration; it is the enabling step for the paper's headline
result, not a performance side quest. A second-order benefit: at large
steps, first-order splitting error would be absorbed into the
calibrated $n$ as discretization bias --- the IMEX scheme keeps the
recovered roughness answering for physics, not for the splitting.

Two remaining qualifications. (i)~The gain \emph{grows with
calibration}: $n=0.015$ is ``smooth,'' and realistic calibrated
urban/vegetated values (0.03--0.1) raise $g n^2$ by
$4$--$40\times$, pushing the binding depth of
Eq.~\eqref{eq:threshold} toward tens of centimeters --- the Manning
maps this plan produces make the IMEX path more valuable, not less.
(ii)~These are frozen-coefficient linear-stability estimates, good to
a factor of $\sim$2; the true numbers fall out for free once
\texttt{TSARKIMEX} runs with \texttt{TSAdapt} reporting its own steps.

\bibliographystyle{unsrt}
\bibliography{pi-briefing-manning-calibration}

\end{document}
